\documentclass[a4paper, 12pt]{article}

% ── Paquetes ──────────────────────────────────────────────────────────────────
\usepackage[margin=2.54cm]{geometry}
\usepackage{multicol}
\usepackage{amsmath, amssymb}
\setlength{\marginparwidth}{2cm}
\usepackage{todonotes}   % para las notas pendientes (\todo{...})
\usepackage[fleqn]{amsmath}

% ── Metadatos ─────────────────────────────────────────────────────────────────
\title{Hoja de fórmulas}
\author{}
\date{}

% ══════════════════════════════════════════════════════════════════════════════
\begin{document}
\maketitle

\begin{multicols}{2}

% ── Relatividad especial ──────────────────────────────────────────────────────
\section*{Relatividad especial}

\begin{align*}
    E          &= \sqrt{m^{2} + |\vec{p}|^{2}}                         \\
    P_{i}^{2}  &= P_{f}^{2}
        \quad \text{(Invariantes de Lorentz)}                           \\
    \vec{p}    &= E\,\frac{\vec{u}}{c^{2}}
\end{align*}

\[
\begin{cases}
    A^{0'} = \gamma(A^{0} - v A^{x}) \\
    A^{x'} = \gamma(A^{x} - v A^{0}) \\
    A^{y'} = A^{y} \;;\; A^{z'} = A^{z}
\end{cases}
\]

% ── Isospin ───────────────────────────────────────────────────────────────────
\section*{Isospin}

$\hat{I}_{1},\, \hat{I}_{2},\, \hat{I}_{3} \in SU(2)$

\begin{align*}
    [\hat{I}_{i},\, \hat{I}_{j}]
        &= i\,\epsilon_{ijk}\,\hat{I}_{k}                              \\
    \hat{I}^{2}
        &= \hat{I}_{1}^{2} + \hat{I}_{2}^{2} + \hat{I}_{3}^{2}       \\
    \hat{I}_{\pm}
        &= \hat{I}_{1} \pm i\,\hat{I}_{2}                             \\
    \mu_{jm}
        &= \mu_{j}                                                     \\
    \langle j', m'|\hat{U}|j, m\rangle
        &= \delta_{jj'}\,\delta_{mm'}\,M_{j}                          \\
    \langle j, m|\hat{U}|j, m\rangle
        &= \langle j, m'|\hat{U}|j, m'\rangle
\end{align*}

Postulado:
\begin{equation*}
    [\hat{H}_{F},\, \hat{I}_{j}] = 0
\end{equation*}

% ── Modelo del quark ──────────────────────────────────────────────────────────
\section*{Modelo del quark}

Se utiliza $SU(3)$. \textbf{Matrices de Gell-Mann:}
\todo[inline]{Agregar las matrices de Gell-Mann (Clase 6)}

Algunas relaciones:
\begin{align*}
    \left[\tfrac{\lambda_{1}}{2},\, \tfrac{\lambda_{2}}{2}\right]
        &= i\,\tfrac{\lambda_{3}}{2}                                   \\
    \left[\tfrac{\lambda_{3}}{2},\, \tfrac{\lambda_{5}}{2}\right]
        &= \tfrac{i}{4}\!\left(\lambda_{3} + \sqrt{3}\,\lambda_{8}\right) \\
    \left[\tfrac{\lambda_{6}}{2},\, \tfrac{\lambda_{7}}{2}\right]
        &= \tfrac{i}{4}\!\left(-\lambda_{3} + \sqrt{3}\,\lambda_{8}\right) \\
    \hat{I}_{\pm}
        &= \frac{\lambda_{1} \pm i\lambda_{2}}{2}                     \\
    \hat{V}_{\pm}
        &= \frac{\lambda_{4} \pm i\lambda_{5}}{2}                     \\
    \hat{U}_{\pm}
        &= \frac{\lambda_{6} \pm i\lambda_{7}}{2}                     \\
    Q   &= I_{3} + \tfrac{1}{2}\,Y
\end{align*}

El sabor debe ser un singlete de $SU(3)_{\text{color}}$.

\subsection*{Antifundamental}
Todas las cargas se invierten y agregan signo.

\subsection*{Momento magnético}
\begin{align*}
    \hat{\vec{\mu}}
        &= \frac{q}{2m}\,\hat{\vec{\sigma}}
            \quad (\hat{\vec{\sigma}}: \text{operador de spin})        \\
    \hat{\mu}_{z,i}
        &= Q_{i}\,\frac{e}{2m_{i}}\,\hat{\sigma}_{3}
            \qquad i = u, d, s
\end{align*}

% ── Dirac y Klein-Gordon ──────────────────────────────────────────────────────
\section*{Dirac y Klein-Gordon}

\begin{align*}
    (\partial^{\mu}\partial_{\mu} + m^{2})\,\phi
        &= 0                                                           \\
    (i\gamma^{\mu}\partial_{\mu} - m)\,\psi
        &= 0                                                           \\
    (\vec{k}\cdot\vec{\sigma})^{-1}
        &= \frac{\vec{k}\cdot\vec{\sigma}}{|\vec{k}|^{2}}             \\
    \{\gamma^{\mu},\, \gamma^{\nu}\}
        &= 2\eta^{\mu\nu}                                              \\
    (\gamma^{\mu})^{\dagger}
        &= \gamma^{0}\gamma^{\mu}\gamma^{0}                           \\
    (\gamma^{0})^{\dagger}
        &= \gamma^{0}                                                  \\
    (\gamma^{i})^{\dagger}
        &= -\gamma^{i}
\end{align*}

\todo[inline]{Agregar las gammas y el operador $p$ (Clase 6)}

\subsection*{Transformaciones de Lorentz}
\begin{align*}
    S(\omega)
        &= e^{-\frac{i}{2}\,\omega_{\alpha\beta}\,\Sigma^{\alpha\beta}} \\
    \Sigma^{\alpha\beta}
        &= \frac{1}{4}\!\left[\gamma^{\alpha},\, \gamma^{\beta}\right]
\end{align*}

\subsection*{Helicidad}
\begin{align*}
    h   &= \vec{s}\cdot\hat{p}
          = \frac{1}{2}\,\frac{\vec{\Sigma}\cdot\vec{p}}{|\vec{p}|}  \\
    [H,\, h] &= 0
\end{align*}

\subsection*{Quiralidad}
\begin{align*}
    \gamma^{5}
        &= i\,\gamma^{0}\gamma^{1}\gamma^{2}\gamma^{3}                \\
    [\gamma^{5},\, S(\Lambda)]
        &= 0                                                           \\
    [H,\, \gamma^{5}]
        &\neq 0                                                        \\
    \gamma^{5}
        &\xrightarrow{m\to 0} 2h                                      \\
    (\gamma^{5})^{\dagger}
        &= \gamma^{5}                                                  \\
    (\gamma^{5})^{2}
        &= 1                                                           \\
    \{\gamma^{5},\, \gamma^{\mu}\}
        &= 0                                                           \\
    \hat{P}_{R}
        &= \tfrac{1}{2}(1 + \gamma^{5})                               \\
    \hat{P}_{L}
        &= \tfrac{1}{2}(1 - \gamma^{5})
\end{align*}

% ── Lagrangiano ───────────────────────────────────────────────────────────────
\section*{Lagrangiano}

\begin{align*}
    \partial_{\mu}\!\left(\frac{\partial\mathcal{L}}
        {\partial(\partial_{\mu}\phi)}\right)
        - \frac{\partial\mathcal{L}}{\partial\phi}
        &= 0                                                           \\
    \frac{\partial(\partial_{\mu}\phi)}{\partial(\partial_{\nu}\phi)}
        &= \delta^{\nu}_{\mu}                                          \\
    \frac{\partial(\partial^{\mu}\phi)}{\partial(\partial_{\nu}\phi)}
        &= \eta^{\mu\nu}                                               \\
    \frac{\partial(\partial_{\mu}A_{\nu})}
         {\partial(\partial_{\alpha}A_{\beta})}
        &= \delta^{\alpha}_{\mu}\,\delta^{\beta}_{\nu}
\end{align*}

\subsection*{Teorema de Noether}

Sea $\mathcal{L} = \mathcal{L}(\phi_{1},\dots,\phi_{N}, \partial_{\mu}\phi_{1},\dots,\partial_{\mu}\phi_{N})$. Ante una transformación infinitesimal tal que $\mathcal{L}' = \mathcal{L} + \partial_{\mu}K^{\mu}$:

\begin{equation*}
    J^{\mu} = \sum_{i}
        \frac{\partial\mathcal{L}}{\partial(\partial_{\mu}\phi_{i})}
        \,\Delta\phi_{i} - K^{\mu},
    \qquad
    \partial_{\mu}J^{\mu} = 0
    \;\text{(on-shell)}
\end{equation*}

\end{multicols}
\end{document}